% Prose rendering of PrimeTensor/Analysis/Near/Laws.lean.
% Fragment: no \documentclass, no preamble, no document environment.

\section{Algebraic laws for nearness and convergence}
\label{Analysis:Near:Laws:sec}

\leanfile{PrimeTensor/Analysis/Near/Laws.lean}

\subsection*{What this file establishes}

\texttt{PrimeTensor/Analysis/Limit.lean} introduced nearness and convergence on
$\name{MulReal}$ but related them to nothing: neither the multiplication of
\texttt{PrimeTensor/Algebra.lean} nor its inversion was known to interact with
them.  This file supplies the four compatibilities, and one consequence.
\[
  \begin{array}{lll}
    \text{nearness} & \text{multiplication} &
      \text{costs one level} \\
    \text{nearness} & \text{inversion} &
      \text{costs nothing} \\
    \text{convergence} & \text{multiplication} &
      \text{costs one level} \\
    \text{convergence} & \text{inversion} &
      \text{costs nothing}
  \end{array}
\]
The consequence is that ratios of convergent sequences converge to the ratio of
the limits.  Together the four are the sequential statement that the carrier
operations are continuous.

Nothing new is estimated.  The two nearness laws transport
$\name{scaleWithin\_mul}$ and $\name{scaleWithin\_inv}$ of
\texttt{PrimeTensor/Algebra.lean} across the quotient, and the two convergence
laws are the anchor-combining pattern used three times already.  What the file
does show, for the first time, is that the deliberately weak nearness relation
of \texttt{PrimeTensor/Analysis/Limit.lean} is strong enough to be transported
at all --- see Remark~\ref{Analysis:Near:Laws:rem:existential}.

\subsection{Nearness respects the operations}

\begin{theorem}[\lean{MulReal.scaleNear\_mul}]
\label{Analysis:Near:Laws:thm:near-mul}
If
\[
  \name{ScaleNear}(\ctor{succ}\,\ell,\, a,\, a')
  \quad\text{and}\quad
  \name{ScaleNear}(\ctor{succ}\,\ell,\, b,\, b'),
\]
then $\name{ScaleNear}(\ell,\, a b,\, a' b')$.
\end{theorem}

\begin{leancode}
theorem scaleNear_mul {level : Depth} {a a' b b' : MulReal}
    (ha : ScaleNear (.succ level) a a')
    (hb : ScaleNear (.succ level) b b') :
    ScaleNear level (a * b) (a' * b') := by
  obtain ⟨as, as', has, has', aAnchor, haTail⟩ := ha
  obtain ⟨bs, bs', hbs, hbs', bAnchor, hbTail⟩ := hb

  refine ⟨MulCauchyStream.mul as bs,
    MulCauchyStream.mul as' bs', ?_, ?_, ?_⟩

  · rw [← has, ← hbs]
    rfl

  · rw [← has', ← hbs']
    rfl
\end{leancode}

\begin{proof}
Each hypothesis supplies a pair of representative streams and an anchor: say
$(\alpha, \alpha')$ with anchor $c_a$ for the first, and $(\beta, \beta')$ with
anchor $c_b$ for the second.  Offer as witnesses for the conclusion the
pointwise products $\alpha\beta$ and $\alpha'\beta'$, using the stream
multiplication of \texttt{PrimeTensor/Algebra.lean}.

The two identifications are $\ctor{rfl}$ after rewriting by the given ones:
$\name{MulReal}$'s multiplication was defined by lifting the stream
multiplication along the quotient, so the class of a product of streams is by
construction the product of their classes, with no lemma needed.

For the tail, take $\name{join}(c_a, c_b)$.  For $n$ at or after it, the two
bounds on $\name{join}$ and transitivity of $\name{AtOrAfter}$, all from
\texttt{PrimeTensor/Completion.lean}, put $n$ in both original tails, so both
factor pairs are $\name{ScaleWithin}(\ctor{succ}\,\ell)$-close at $n$; and
$\name{scaleWithin\_mul}$ of \texttt{PrimeTensor/Algebra.lean} combines them
into $\name{ScaleWithin}(\ell)$-closeness of the products.
\end{proof}

\begin{theorem}[\lean{MulReal.scaleNear\_inv}]
\label{Analysis:Near:Laws:thm:near-inv}
$\name{ScaleNear}(\ell, a, b)$ implies $\name{ScaleNear}(\ell, a^{-1}, b^{-1})$.
\end{theorem}

\begin{proof}
Invert the two witnessing streams pointwise; the identifications are again
$\ctor{rfl}$, because inversion on $\name{MulReal}$ was also defined by lifting.
Keep the anchor and the level, and apply $\name{scaleWithin\_inv}$ termwise.
\end{proof}

\begin{remark}[The existential form is what makes this work]
\label{Analysis:Near:Laws:rem:existential}
$\name{ScaleNear}$ was defined in
\texttt{PrimeTensor/Analysis/Limit.lean} by quantifying \emph{existentially}
over representatives, and that file's design note observed that the choice
makes the relation weaker and easier to establish.  Here it is what makes it
transportable.  Both proofs above proceed by \emph{constructing} witnesses for
the conclusion out of the witnesses supplied by the hypotheses, and the
construction is available because the operations act pointwise on streams and
because $\name{MulReal}$'s operations were defined as their lifts --- so the
required identifications hold definitionally rather than by an argument.

Under a universally quantified nearness the same statements would have to be
proved for representatives one does not get to choose, and each proof would
have to bridge from an arbitrary representative to a convenient one, spending
levels.  So the two theorems above would not have the level costs they do.
\end{remark}

\subsection{Convergence respects the operations}

\begin{theorem}[\lean{MulReal.converges\_mul} and \lean{converges\_inv}]
\label{Analysis:Near:Laws:thm:converges}
If $a \to x$ and $b \to y$ then $a_n b_n \to xy$; and if $a \to x$ then
$a_n^{-1} \to x^{-1}$.
\end{theorem}

\begin{leancode}
theorem converges_mul {a b : Seq} {x y : MulReal}
    (ha : ConvergesTo a x)
    (hb : ConvergesTo b y) :
    ConvergesTo (fun n => a n * b n) (x * y) := by
  intro level
  obtain ⟨aAnchor, haTail⟩ := ha (.succ level)
  obtain ⟨bAnchor, hbTail⟩ := hb (.succ level)
\end{leancode}

\begin{proof}
For the product, fix a level $\ell$, instantiate both hypotheses at
$\ctor{succ}\,\ell$, take $\name{join}$ of the two anchors, and for $n$ beyond
it apply Theorem~\ref{Analysis:Near:Laws:thm:near-mul} to the two nearness
facts.  For the inverse, instantiate at $\ell$, keep the anchor, and apply
Theorem~\ref{Analysis:Near:Laws:thm:near-inv} pointwise.
\end{proof}

\begin{corollary}[\lean{MulReal.converges\_ratio}]
\label{Analysis:Near:Laws:cor:ratio}
If $a \to x$ and $b \to y$ then
$\name{ratio}(a_n, b_n) \to \name{ratio}(x,y)$.
\end{corollary}

\begin{proof}
Unfold $\name{ratio}$ to $a_n \cdot b_n^{-1}$ and compose the two theorems.
\end{proof}

\begin{designnote}[One motif, five layers]\label{Analysis:Near:Laws:note:motif}
The pair ``multiplication costs one level, inversion costs none, and combining
two hypotheses needs a $\name{join}$ of anchors'' has now been proved five
times, at every layer the development has:
\begin{center}
\small
\begin{tabular}{@{}lll@{}}
\hline
Layer & Multiplication & Inversion \\
\hline
$\name{MulRat}$, single pair
  & $\name{scaleWithin\_mul}$ & $\name{scaleWithin\_inv}$ \\
Cauchy streams
  & $\name{MulCauchyStream.mul}$ & $\name{MulCauchyStream.inv}$ \\
asymptotic equivalence
  & $\name{MulAsymptotic.mul}$ & $\name{MulAsymptotic.inv}$ \\
$\name{MulReal}$, nearness
  & Thm.~\ref{Analysis:Near:Laws:thm:near-mul}
  & Thm.~\ref{Analysis:Near:Laws:thm:near-inv} \\
$\name{MulReal}$, convergence
  & \multicolumn{2}{l}{Thm.~\ref{Analysis:Near:Laws:thm:converges}} \\
\hline
\end{tabular}
\end{center}
The first three are in \texttt{PrimeTensor/Algebra.lean} and the last two here.
The repetition is not redundancy: each layer's statement is about different
objects, and only the last two could be stated at all before this file.  But
the \emph{reason} for the level cost is the same at every layer, and it is the
one identified in \texttt{PrimeTensor/Algebra.lean} --- multiplication combines
two independent errors and their ratios multiply, whereas inversion merely
reverses orientation and multiplicative distance does not notice.  The cost is
a property of the operation, not of the layer, so it propagates unchanged all
the way up.
\end{designnote}

\begin{remark}[This is continuity, sequentially]\label{Analysis:Near:Laws:rem:continuity}
Theorem~\ref{Analysis:Near:Laws:thm:converges} is the statement that
multiplication and inversion are continuous, in the sequential form available
here: they carry convergent sequences to convergent sequences with the expected
limits.  Since \texttt{PrimeTensor/Analysis/Limit.lean} defined convergence for
sequences and not through a filter or a topology, this is the strongest form
the statement can take in the present vocabulary.  In particular nothing here
asserts joint continuity in any uniform sense, and nothing asserts that the
limit is unique --- that gap, open since
\texttt{PrimeTensor/Structure.lean}, is untouched, so ``the'' limit still may
not be spoken of.
\end{remark}

\begin{remark}[What is not proved]\label{Analysis:Near:Laws:rem:missing}
Nearness is still not transitive at any level, for the reason
\texttt{PrimeTensor/Analysis/Limit.lean} gave, and nothing here changes that.
There is no compatibility of nearness or convergence with $\name{ofRat}$, no
statement that a convergent sequence is Cauchy in any sense --- the
vocabulary for that is still missing --- and no interaction with the
differential notions of
\texttt{PrimeTensor/Analysis/Derivative.lean}: in particular
$\name{NegligibleRelative}$ is not shown to be preserved by products, which is
the fact a product rule would need.
\end{remark}

\subsection*{Summary of the correspondence}

\begin{center}
\small
\begin{tabular}{@{}lll@{}}
\hline
Lean declaration & Kind & Here \\
\hline
\lean{MulReal.scaleNear\_mul}   & theorem & Thm.~\ref{Analysis:Near:Laws:thm:near-mul} \\
\lean{MulReal.scaleNear\_inv}   & theorem & Thm.~\ref{Analysis:Near:Laws:thm:near-inv} \\
\lean{MulReal.converges\_mul}   & theorem & Thm.~\ref{Analysis:Near:Laws:thm:converges} \\
\lean{MulReal.converges\_inv}   & theorem & Thm.~\ref{Analysis:Near:Laws:thm:converges} \\
\lean{MulReal.converges\_ratio} & theorem & Cor.~\ref{Analysis:Near:Laws:cor:ratio} \\
\hline
\end{tabular}
\end{center}

\noindent
Every declaration of \texttt{PrimeTensor/Analysis/Near/Laws.lean} appears
above.  The file contains no \lean{sorry}, no axiom, and no unproved named
proposition.  Design note~\ref{Analysis:Near:Laws:note:motif} and
Remarks~\ref{Analysis:Near:Laws:rem:existential},
\ref{Analysis:Near:Laws:rem:continuity}
and~\ref{Analysis:Near:Laws:rem:missing} are stated for the reader and are not
declarations of the file.
